% ------------------------------------------------------------
% Sprache, Zeichencodierung und Anführungen
% ------------------------------------------------------------
\usepackage[ngerman]{babel}
\usepackage[utf8]{inputenc}
\usepackage[T1]{fontenc}
\usepackage{csquotes}

% ------------------------------------------------------------
% Seitenlayout: DIN A4, 2 cm Ränder
% ------------------------------------------------------------
\usepackage{geometry}
\geometry{
  top=2cm,
  bottom=2cm,
  left=2cm,
  right=2cm
}

% ------------------------------------------------------------
% Zeilenabstand und Feintypografie
% ------------------------------------------------------------
\usepackage{setspace}
\onehalfspacing
\usepackage{microtype}

% ------------------------------------------------------------
% Schrift: TeX Gyre Heros als Arial-ähnliche Schrift
% ------------------------------------------------------------
\usepackage{tgheros}
\renewcommand{\familydefault}{\sfdefault}

% ------------------------------------------------------------
% Überschriftenformatierung
% ------------------------------------------------------------
\usepackage{titlesec}

\titleformat{\section}
  {\bfseries\fontsize{12pt}{14pt}\selectfont}
  {\thesection}
  {1em}
  {}

\titleformat{\subsection}
  {\bfseries\fontsize{12pt}{14pt}\selectfont}
  {\thesubsection}
  {1em}
  {}

\titleformat{\subsubsection}
  {\bfseries\fontsize{12pt}{14pt}\selectfont}
  {\thesubsubsection}
  {1em}
  {}

\setcounter{secnumdepth}{3}
\setcounter{tocdepth}{3}

% ------------------------------------------------------------
% Absätze und Fußnoten
% ------------------------------------------------------------
\setlength{\parskip}{6pt}
\setlength{\parindent}{0pt}

\makeatletter
\renewcommand\footnotesize{\@setfontsize\footnotesize{10pt}{12pt}}
\makeatother

% ------------------------------------------------------------
% Seitenzahlen: zentriert unten, keine Kopfzeile
% ------------------------------------------------------------
\usepackage{fancyhdr}
\pagestyle{fancy}
\fancyhf{}
\fancyfoot[C]{\thepage}
\renewcommand{\headrulewidth}{0pt}

% ------------------------------------------------------------
% Grafiken, Tabellen und codebasierte Abbildungen
% ------------------------------------------------------------
\usepackage{graphicx}
\usepackage{xcolor}
\ifbeginnerstatements
  \usepackage{tcolorbox}
\fi
\usepackage{tikz}
\usetikzlibrary{
  arrows.meta,
  backgrounds,
  calc,
  fit,
  matrix,
  positioning,
  shapes.geometric
}

% Einheitliche Farben der Fallstudiengrafiken
\definecolor{caseBlue}{HTML}{174A7E}
\definecolor{caseGreen}{HTML}{5A8F3D}
\definecolor{caseLight}{HTML}{F4F8FB}
\definecolor{caseLine}{HTML}{A9B4BF}

% Optionale, zurückhaltende Erklärungen für technische Neueinsteiger
\ifbeginnerstatements
  \definecolor{beginnerBackground}{HTML}{F3F3F3}
  \definecolor{beginnerBorder}{HTML}{B8B8B8}

  \newtcolorbox{beginnerstatement}[1][Kurz erklärt]{
    colback=beginnerBackground,
    colframe=beginnerBorder,
    colbacktitle=beginnerBackground,
    coltitle=black,
    coltext=black,
    boxrule=0.6pt,
    arc=2mm,
    outer arc=2mm,
    left=2.5mm,
    right=2.5mm,
    top=2mm,
    bottom=2mm,
    before skip=10pt plus 2pt minus 2pt,
    after skip=10pt plus 2pt minus 2pt,
    fonttitle=\bfseries,
    fontupper=\small,
    title={#1}
  }
\fi

\tikzset{
  casePrimary/.style={
    draw=caseBlue,
    fill=caseBlue,
    text=white,
    rounded corners=3pt,
    align=center,
    font=\sffamily\bfseries\scriptsize,
    minimum height=8mm,
    inner sep=4pt
  },
  caseBox/.style={
    draw=caseBlue,
    fill=white,
    text=caseBlue,
    rounded corners=3pt,
    align=center,
    font=\sffamily\scriptsize,
    minimum height=8mm,
    inner sep=4pt
  },
  caseOptional/.style={
    draw=caseGreen,
    fill=caseGreen!10,
    text=caseGreen!55!black,
    rounded corners=3pt,
    align=center,
    font=\sffamily\scriptsize,
    minimum height=8mm,
    inner sep=4pt
  },
  caseGroup/.style={
    draw=caseLine,
    fill=caseLight,
    rounded corners=5pt,
    inner sep=6pt
  },
  caseFlow/.style={
    -{Latex[length=2.6mm,width=1.8mm]},
    draw=caseBlue,
    line width=0.9pt
  },
  caseSecondaryFlow/.style={
    -{Latex[length=2.4mm,width=1.7mm]},
    draw=caseLine,
    line width=0.8pt,
    dashed
  }
}

\usepackage{booktabs}
\usepackage{array}
\usepackage{tabularx}
\usepackage{float}
\usepackage{threeparttable}

% ------------------------------------------------------------
% Hyperlinks
% ------------------------------------------------------------
\usepackage[hidelinks]{hyperref}
\usepackage{url}

% ------------------------------------------------------------
% Literaturverwaltung mit biblatex-apa
% ------------------------------------------------------------
\usepackage[
  style=apa,
  backend=biber,
  sorting=nyt,
  giveninits=true,
  maxcitenames=2,
  maxbibnames=20,
  doi=true,
  url=true,
  isbn=true
]{biblatex}

\DeclareLanguageMapping{ngerman}{ngerman-apa}

\DefineBibliographyStrings{ngerman}{
  nodate = {o\adddot\addspace D\adddot}
}

\addbibresource{references.bib}
